\documentclass[11pt,letterpaper]{article}
\usepackage[spanish,es-nodecimaldot]{babel}
\usepackage[utf8]{inputenc}
\usepackage[T1]{fontenc}
\usepackage{lmodern,microtype}
\usepackage[margin=25mm,inner=30mm,top=30mm]{geometry}
\usepackage{setspace}
\setstretch{1.3}
\setlength{\parindent}{0pt}
\setlength{\parskip}{5pt}
\usepackage{xcolor,graphicx,float,booktabs,caption}
\usepackage{fancyhdr}
\usepackage[colorlinks=true,linkcolor=black,urlcolor=blue]{hyperref}
\usepackage{xurl}
\usepackage{listings}
\lstdefinelanguage{MySQLLab}{
  language=SQL,
  morekeywords={AUTO_INCREMENT,DECIMAL,ENUM,SHOW,SET,SQL_SAFE_UPDATES},
  sensitive=false
}
\lstdefinestyle{sqlusm}{
 language=MySQLLab,basicstyle=\ttfamily\footnotesize,
 keywordstyle=\color{azulusm}\bfseries,
 commentstyle=\color{gray},stringstyle=\color{teal!65!black},
 breaklines=true,breakatwhitespace=false,
 backgroundcolor=\color{sqlbg},frame=single,rulecolor=\color{black!25},
 columns=fullflexible,keepspaces=true,showstringspaces=false,
 xleftmargin=3pt,xrightmargin=3pt,
 literate={á}{{\'a}}1 {é}{{\'e}}1 {í}{{\'i}}1 {ó}{{\'o}}1 {ú}{{\'u}}1
 {Á}{{\'A}}1 {É}{{\'E}}1 {Í}{{\'I}}1 {Ó}{{\'O}}1 {Ú}{{\'U}}1
 {ñ}{{\~n}}1 {Ñ}{{\~N}}1 {ü}{{\"u}}1
}

\usepackage{ifthen}
\usepackage{needspace}
\usepackage{enumitem}
\usepackage[section]{placeins}
\definecolor{azulusm}{HTML}{174B79}
\definecolor{sqlbg}{HTML}{F4F6F9}
\lstset{style=sqlusm}
\newcommand{\theCampus}{Campus Vitacura}
\newcommand{\theAssignment}{Laboratorio N.º 4 — Bases de Datos y SQL}
\newcommand{\therevision}{ICN292 -- USM}
\pagestyle{fancy}
\fancyhf{}
\fancyhead[L]{\IfFileExists{Logo_UTFSM.png}{\includegraphics[height=0.9cm]{Logo_UTFSM.png}}{\small\bfseries Universidad Técnica Federico Santa María}}
\fancyfoot[L]{\scriptsize Laboratorio N.º 4}
\fancyfoot[C]{\thepage}
\fancyfoot[R]{\scriptsize ICN292 -- USM}
\setlength{\headheight}{36pt}
\renewcommand{\headrulewidth}{0.4pt}
\renewcommand{\footrulewidth}{0.4pt}
\captionsetup{font=small,labelfont=bf}
\newcommand{\evidencia}[3][0.95]{
  \begin{figure}[H]\centering
  \includegraphics[width=#1\linewidth,height=0.27\textheight,keepaspectratio]{capturas/#2.png}
  \caption{#3}\end{figure}
}
\newcommand{\capturapendiente}{
\begin{center}\fcolorbox{black!40}{sqlbg}{\parbox{.85\linewidth}{\centering\textbf{Captura pendiente:} adjuntar la evidencia del \texttt{UPDATE} ejecutado correctamente y el mensaje de filas afectadas de MySQL Workbench.}}\end{center}}
\begin{document}
\begin{titlepage}
\thispagestyle{empty}
\begin{center}
\vspace*{1cm}
\IfFileExists{header utfsm.png}{\includegraphics[width=0.60\textwidth]{header utfsm.png}}{
{\Large\bfseries UNIVERSIDAD TÉCNICA FEDERICO SANTA MARÍA}\\[0.25cm]
{\large Departamento de Industrias — Campus Vitacura}}
\vspace{1.6cm}
\par\noindent\rule{\linewidth}{0.8pt}\par\vspace{0.2cm}
{\LARGE\bfseries Informe de Laboratorio N.º 4}\\[0.35cm]
{\Large Bases de Datos y SQL — Austral Editores}\\[0.35cm]
{\large Sistemas de Información para la Gestión (ICN292)}
\vspace{0.2cm}\par\noindent\rule{\linewidth}{0.8pt}\par
\vspace{1.5cm}
{\Large\bfseries Sergio Martínez}\\[0.5cm]
{\large Semestre 2026/02}\\[0.35cm]
{\large Departamento de Industrias}\\[0.5cm]
\vspace{0.7cm}
\textbf{Repositorio público GitHub:}\\[0.2cm]
\href{https://github.com/Sergi9787/ICN292-Lab4-Martinez-Sergio}{\texttt{https://github.com/Sergi9787/ICN292-Lab4-Martinez-Sergio}}\\
\vfill
{\large Noviembre 2026}\\[0.35cm]
{\large\theCampus}
\end{center}
\end{titlepage}
\section*{Entregable 2: Consultas SQL}
\addcontentsline{toc}{section}{Entregable 2: Consultas SQL}
\subsection*{Consulta 1. Autores cuyo país no es Chile}
\begin{lstlisting}
SELECT nombre, genero
FROM autores
WHERE pais <> 'Chile' OR pais IS NULL
ORDER BY nombre ASC;
\end{lstlisting}
\textbf{Explicación solicitada:} \texttt{pais <> 'Chile'} devuelve los países conocidos distintos de Chile; el valor \texttt{NULL} no satisface esa comparación porque produce un valor lógico desconocido, por lo que se agrega \texttt{OR pais IS NULL}.
\evidencia[0.7]{c01}{Consulta 1: autores seleccionados, incluyendo el país desconocido.}
\FloatBarrier

\subsection*{Consulta 2. Libros de 190 a 220 páginas}
\begin{lstlisting}
SELECT titulo, paginas
FROM libros
WHERE paginas BETWEEN 190 AND 220
ORDER BY paginas ASC, titulo ASC;
\end{lstlisting}
\evidencia[0.77]{c02}{Consulta 2: siete libros dentro del intervalo de páginas.}
\FloatBarrier

\subsection*{Consulta 3. Indicadores del catálogo original}
\begin{lstlisting}
SELECT COUNT(*) AS total_libros,
       SUM(ejemplares_vendidos) AS total_ejemplares,
       AVG(paginas) AS promedio_paginas,
       MIN(anio) AS anio_mas_antiguo,
       MAX(ejemplares_vendidos) AS mayor_venta_libro
FROM libros;
\end{lstlisting}
\evidencia[1.0]{c03}{Consulta 3: indicadores consolidados de la carga oficial.}
\FloatBarrier

\subsection*{Consulta 4. Géneros distintos}
\begin{lstlisting}
SELECT DISTINCT CAST(genero AS CHAR) AS genero
FROM autores
ORDER BY genero ASC;
\end{lstlisting}
\evidencia[0.53]{c04}{Consulta 4: géneros sin duplicados, en orden alfabético.}
\FloatBarrier

\subsection*{Consulta 5. Autores con más de un libro en 2024}
\begin{lstlisting}
SELECT a.nombre, COUNT(*) AS libros_2024
FROM autores AS a
JOIN libros AS l ON l.id_autor = a.id_autor
WHERE l.anio = 2024
GROUP BY a.id_autor, a.nombre
HAVING COUNT(*) > 1
ORDER BY a.nombre ASC;
\end{lstlisting}
\evidencia[0.65]{c05}{Consulta 5: autores con más de una publicación durante 2024.}
\FloatBarrier

\subsection*{Consulta 6. Tres libros más vendidos de 2024}
\begin{lstlisting}
SELECT titulo, ejemplares_vendidos
FROM libros
WHERE anio = 2024
  AND ejemplares_vendidos > 2000
ORDER BY ejemplares_vendidos DESC, titulo ASC
LIMIT 3;
\end{lstlisting}
\evidencia[0.72]{c06}{Consulta 6: tres libros de 2024 con más de 2.000 ejemplares.}
\FloatBarrier

\subsection*{Consulta 7. Inserción y eliminación de Papel Mojado}
\begin{lstlisting}
INSERT INTO libros (id_autor, titulo, paginas,
                   ejemplares_vendidos, anio)
VALUES (3, 'Papel Mojado', 260, 50, 2024);

SELECT * FROM libros
WHERE titulo = 'Papel Mojado';
\end{lstlisting}
\begin{lstlisting}
-- Para ejecutar el borrado por titulo en MySQL Workbench,
-- es necesario que la sesión lo permita (Safe Updates).
SET SQL_SAFE_UPDATES = 0;
DELETE FROM libros WHERE titulo = 'Papel Mojado';
SET SQL_SAFE_UPDATES = 1;

SELECT * FROM libros
WHERE titulo = 'Papel Mojado';
\end{lstlisting}
\evidencia[1.0]{c07_antes}{Consulta 7: libro insertado, antes de eliminarlo.}
\evidencia[1.0]{c07_despues}{Consulta 7: SELECT posterior a la eliminación, sin resultados.}
\FloatBarrier

\subsection*{Consulta 8. Aumento de 100 ejemplares en Cielo de Enero}
\begin{lstlisting}
SELECT titulo, ejemplares_vendidos
FROM libros
WHERE titulo = 'Cielo de Enero';
\end{lstlisting}
\begin{lstlisting}
UPDATE libros
SET ejemplares_vendidos = ejemplares_vendidos + 100
WHERE id_libro = 1;
\end{lstlisting}
\begin{lstlisting}
SELECT titulo, ejemplares_vendidos
FROM libros
WHERE titulo = 'Cielo de Enero';
\end{lstlisting}
\textbf{Evidencia adicional requerida por el enunciado:} captura de la ejecución exitosa del \texttt{UPDATE}, todavía no proporcionada.\par\capturapendiente
\evidencia[0.75]{c08_antes}{Consulta 8: SELECT previo, con 1.500 ejemplares.}
\evidencia[0.75]{c08_despues}{Consulta 8: SELECT posterior, con 1.600 ejemplares.}
\FloatBarrier

\subsection*{Consulta 9. Colecciones y sus libros}
\begin{lstlisting}
-- Datos de prueba incorporados para la consulta 9.
INSERT INTO colecciones (titulo, fecha_publicacion) VALUES
('Colección Austral Contemporánea', '2024-03-14'),
('Colección Letras y Horizontes', '2023-08-09');

UPDATE libros SET id_coleccion = 1
WHERE id_libro IN (1, 3, 5, 8, 9, 11);
UPDATE libros SET id_coleccion = 2
WHERE id_libro IN (2, 4, 6, 7, 10, 12);
\end{lstlisting}
\begin{lstlisting}
SELECT l.titulo AS libro,
       c.titulo AS coleccion,
       c.fecha_publicacion
FROM libros AS l
JOIN colecciones AS c
    ON l.id_coleccion = c.id_coleccion
ORDER BY l.id_libro;
\end{lstlisting}
\evidencia[1.0]{c09}{Consulta 9: doce libros vinculados a sus colecciones.}
\FloatBarrier

\subsection*{Consulta 10. Pagos a los autores}
\begin{lstlisting}
-- Datos de prueba incorporados para la consulta 10.
INSERT INTO pagos (id_autor, fecha_pago, monto) VALUES
(1, '2026-10-01', 150000.00),
(1, '2026-10-15',  95000.50),
(3, '2026-10-19',  80000.00),
(4, '2026-10-27', 120000.00);
\end{lstlisting}
\begin{lstlisting}
SELECT a.nombre AS autor, p.fecha_pago, p.monto
FROM pagos AS p
JOIN autores AS a ON p.id_autor = a.id_autor
ORDER BY a.nombre ASC, p.fecha_pago ASC;
\end{lstlisting}
\begin{lstlisting}
SELECT a.nombre AS autor,
       SUM(p.monto) AS total_pagado
FROM autores AS a
JOIN pagos AS p ON a.id_autor = p.id_autor
GROUP BY a.id_autor, a.nombre
ORDER BY a.nombre ASC;
\end{lstlisting}
\evidencia[0.85]{c10_detalle}{Consulta 10, primera sentencia: detalle de los cuatro pagos.}
\evidencia[0.8]{c10_total}{Consulta 10, segunda sentencia: suma de pagos por autor.}
\FloatBarrier

\end{document}
